\subsection{张量与不变形式}

设 G 是李群。我们把 G 看作通过左（分别为右）平移作用于自身。设 \(\lambda\) 是对同构为 \(C^{\omega}\) 类的向量函子。则 \(\lambda(\mathrm{TG})\) 是解析的左（分别为右）向量 G-丛（§1，第8号，命题16的推论）。映射 \((g, u) \mapsto gu\)（分别为 ug）从 \(G \times \lambda(\mathrm{L}(G))\) 到 \(\lambda(\mathrm{TG})\) 上是向量 G-丛的同构 \(\phi\)（分别为 \(\psi\)）（§7，第8号，命题17的推论2）。\(\lambda(\mathrm{TG})\) 的每个 G-不变截面都是解析的，并由其在 e 处的值确定（§1，第8号，命题17的推论1）。这样的截面称为左（分别为右）不变。设 \(\sigma\) 是 \(\lambda(\mathrm{TG})\) 的左不变截面；变换 \(\sigma'\) 是 \(\sigma\) 在右平移 \(\delta(g)\) 下的变换，定义为 \(\sigma'(\delta(g)h) = \lambda(\mathrm{T}_{h}(\delta(g)))\sigma(h)\)，对所有 \(h \in G\)；它也是左不变的；它也由 \(\sigma\) 经 \(\gamma(g) \circ \delta(g) = \operatorname{Int}(g)\) 导出，因而

\[
\sigma^ {\prime} (e) = \lambda (\mathrm{Ad} g). \sigma (e).\tag{26}
\]

类似地，设 \(\tau\) 是 \(\lambda(\mathrm{TG})\) 的右不变截面；变换 \(\tau'\) 是 \(\tau\) 在左平移 \(\gamma(g)\) 下的变换，也具有右变性，并且

\[
\tau^ {\prime} (e) = \lambda (\mathrm{Ad} g). \tau (e).\tag{27}
\]

现在考虑 G × G 按下式在左侧作用于 G：

\[
\left(\left(g, g ^ {\prime}\right), g ^ {\prime \prime}\right) \mapsto g g ^ {\prime \prime} g ^ {\prime - 1}.
\]

于是 G 是 \(G \times G\) 的左李齐性空间（§ 1，第6号，例）。因此 \(\lambda(\mathrm{TG})\) 是解析的左向量 \((\mathrm{G} \times \mathrm{G})\)-丛。若 \(\lambda(\mathrm{TG})\) 的一个截面在 \(\mathrm{G} \times \mathrm{G}\) 对 \(\lambda(\mathrm{TG})\) 的作用下不变，或者说在左、右平移下都不变，则称其为双不变截面。设 \(\lambda(\mathrm{L}(\mathrm{G}))_{0}\) 是 \(\lambda(\mathrm{L}(\mathrm{G}))\) 中在 \(\lambda(\mathrm{Ad}(\mathrm{G}))\) 下不变的元素集合。对任意 \(u \in \lambda(\mathrm{L}(\mathrm{G}))_{0}\)，设 \(\sigma_{u}\) 是从 G 到 \(\lambda(\mathrm{TG})\) 的映射，由 \(\sigma_{u}(g) = gu = ug\) 定义。则 \(u \mapsto \sigma_{u}\) 是从 \(\lambda(\mathrm{L}(\mathrm{G}))_{0}\) 到 \(\lambda(\mathrm{TG})\) 的双不变截面集合的双射（§ 1，第8号，命题17的推论1）。

命题 50. 设 G 是李群（当 K 的特征 \textgreater0 时假定其有限维）。设 E 是 \(T_{e}(G)\) 上 k 阶连续交替多重线性形式的向量空间。对任意 \(u \in E\)，设 \(\omega^{u}\) 是 G 上满足如下条件的 k 阶微分形式：\((\omega^{u})_{g}\) 是 \(T_{e}(G)\) 上由平移 \(h \mapsto gh\)（分别为 \(h \mapsto hg\)）从 u 导出的多重线性形式。则 \(\omega^{u}\) 在 G 上解析且左（分别为右）不变。映射 \(u \mapsto \omega^{u}\) 是 E 到 G 上 k 阶左（分别为右）不变微分形式向量空间的同构。

这是上文所述内容的特例。

设 \(\mathbf{F}\) 是完备可赋范空间。若把 \(\mathbf{G}\) 上取值于 \(\mathbf{K}\) 的微分形式替换为 \(\mathbf{G}\) 上取值于 \(\mathbf{F}\) 的微分形式，命题50仍成立。对于每个连续线性映射 \(u\)，它从 \(\mathrm{T}_e(\mathrm{G})\) 映到 \(\mathbf{F}\)，都存在一个 1 阶微分形式 \(\omega^u\)，定义在 \(\mathbf{G}\) 上且取值于 \(\mathbf{F}\)，使得 \((\omega^{u})_{g} = u \circ \mathrm{T}_{g}(\gamma(g)^{-1})\)。特别地，取 \(\mathrm{F} = \mathrm{T}_e(\mathrm{G})\) 及 \(u = \mathrm{Id}_{\mathrm{T_e(G)}}\)。于是得到 G 上的微分形式 \(\omega\)，满足 \(\omega_{g} = \mathrm{T}_{g}(\gamma (g^{-1}))\)；这个微分形式左不变且解析；称为 G 的左典范微分形式。满足 \(\omega_{g}(t) = g^{-1}t\)，对所有 \(t\in \mathrm{T}_g(\mathrm{G})\)。

若 F 再次为任意完备可赋范空间且 \(u \in \mathcal{L}(\mathrm{T}_{e}(\mathrm{G}), \mathrm{F})\)，则 \(\omega^{u} = u \circ \omega\)。特别地（取 F = K），映射 \(v \mapsto v \circ \omega\) 是从 \(\mathrm{T}_{e}(\mathrm{G})\) 的对偶到取值于 K 且在 G 下左不变的 1 阶微分形式向量空间的线性双射。

类似地，G 上的微分形式 \(\omega'\) 满足 \(\omega_g' = \mathrm{T}_g(\delta(g))\)，称为 G 的右典范微分形式。它具有与 \(\omega\) 类似的性质，留给读者陈述。G 到 G 的映射 \(g \mapsto g^{-1}\) 将 \(\omega\) 变换为 \(\omega'\)。

